% =====================================================================
% 4. Algoritmos e implementação
% Os trechos de código são cópias exatas de codigo/*.sci e codigo/main.sce.
% =====================================================================
\section{Algoritmos e implementação}
\label{sec:algoritmos}

\subsection{Entrada, saída e critério de parada}

O programa foi escrito em Scilab e está dividido em seis arquivos, listados no Apêndice~\ref{ap:codigo}. A entrada de dados fica no início do \texttt{main.sce}. Ela define os componentes do circuito, a amplitude da fonte, os intervalos de simulação e os passos de cada ensaio:
\begin{codigo}
L = 10e-3;          // indutância (H)
C = 1e-6;           // capacitância (F)
V0 = 5;             // amplitude da fonte (V)
R = 20;             // resistência do circuito principal (ohm), Q = 5
tf_conv = 5e-3;     // intervalo dos ensaios de convergência (s)
hs_conv = [8e-6, 4e-6, 2e-6, 1e-6, 5e-7, 2.5e-7];
\end{codigo}
A saída tem duas partes. O arquivo \texttt{codigo/resultados/metricas\_scilab.json} guarda todos os números das tabelas. A pasta \texttt{figuras/} recebe os gráficos em PDF.

A integração tem dois critérios de parada. O normal é chegar ao instante final \(t_f\). O outro detecta divergência: a integração para se algum valor do estado deixa de ser finito ou passa de um limite muito acima das grandezas físicas do circuito. Esse limite é \(10^{6}\) nos ensaios comuns e \(10^{3}\) no ensaio de estabilidade. O segundo critério separa uma solução que explodiu de uma simples oscilação.

\subsection{Pseudocódigo}

O Algoritmo~\ref{alg:rk4} resume a integração com o RK4. Para o Euler e o Heun, muda só o bloco que calcula o novo valor de \(y\): o Euler faz uma avaliação de \(f\) por passo, e o Heun faz duas.

\begin{algoritmo}[htbp]
\centering
\caption{Runge-Kutta de quarta ordem com passo fixo e parada por divergência.}
\label{alg:rk4}
\begin{minipage}{0.94\textwidth}
\setstretch{1.05}
\noindent\textbf{Entrada:} \(f\), \(y_0\), \(t_0\), \(t_f\), \(h\), limite de divergência.\\[0.3em]
\(n\leftarrow\mathrm{round}((t_f-t_0)/h)\);\quad \(h\leftarrow(t_f-t_0)/n\);\quad \(y\leftarrow y_0\);\quad \(N_f\leftarrow 0\).\\[0.3em]
\textbf{para} \(k=0,1,\ldots,n-1\) \textbf{faça}\\[0.2em]
\hspace*{1.2em}\(t\leftarrow t_0+kh\)\\
\hspace*{1.2em}\(K_1\leftarrow h\,f(t,y)\)\\
\hspace*{1.2em}\(K_2\leftarrow h\,f(t+h/2,\; y+K_1/2)\)\\
\hspace*{1.2em}\(K_3\leftarrow h\,f(t+h/2,\; y+K_2/2)\)\\
\hspace*{1.2em}\(K_4\leftarrow h\,f(t+h,\; y+K_3)\)\\
\hspace*{1.2em}\(y\leftarrow y+(K_1+2K_2+2K_3+K_4)/6\)\\
\hspace*{1.2em}\(N_f\leftarrow N_f+4\)\\
\hspace*{1.2em}\textbf{se} \(y\) não é finito ou \(\max|y|>\mathrm{limite}\) \textbf{então} marcar divergência e sair\\[0.3em]
\textbf{Saída:} pares \((t_k,y_k)\), passo usado e número \(N_f\) de avaliações de \(f\).
\end{minipage}
\fonte{Elaborado pelos autores (2026), com base em \textcite[seção~10.3.2]{aquino2020}.}
\end{algoritmo}

\subsection{O código, função por função}

\paragraph{Modelo do circuito.} A função \texttt{rlc\_campo} calcula o lado direito \(f(t,y)\) da Equação~\eqref{eq:sistema}. A primeira linha aplica a lei das malhas e dá \(di/dt\). A segunda aplica a relação do capacitor e dá \(dv_C/dt\). A terceira junta as duas derivadas num vetor coluna. A struct \texttt{p} guarda \(R\), \(L\), \(C\) e os dados da fonte.
\begin{codigo}
function dy = rlc_campo(t, y, p)
    di = (fonte_vs(t, p) - p.R * y(1) - y(2)) / p.L;
    dvc = y(1) / p.C;
    dy = [di; dvc];
endfunction
\end{codigo}

\paragraph{Os três métodos.} Cada método é uma função que avança um único passo e devolve o novo estado \texttt{y1} e o número \texttt{nf} de avaliações de \(f\) que gastou. A forma segue o livro da disciplina \cite[cap.~10]{aquino2020}. O Euler aplica a Equação~\eqref{eq:euler} numa linha só:
\begin{codigo}
y1 = y + h * rlc_campo(t, y, p);
\end{codigo}
O Heun segue a Equação~\eqref{eq:heun}. A variável \texttt{f1} guarda a inclinação no início do passo, para que ela não seja calculada duas vezes. \texttt{K1} é a previsão de Euler. A última linha corrige o passo com a média das duas inclinações:
\begin{codigo}
f1 = rlc_campo(t, y, p);
K1 = y + h * f1;
y1 = y + h * (f1 + rlc_campo(t + h, K1, p)) / 2;
\end{codigo}
O RK4 segue a Equação~\eqref{eq:rk4}. Cada \texttt{K} já vem multiplicado por \(h\), como em \textcite[seção~10.3.2]{aquino2020}. \texttt{K2} e \texttt{K3} usam o meio do passo, \texttt{K4} usa o fim, e a última linha faz a média com pesos 1, 2, 2 e 1:
\begin{codigo}
K1 = h * rlc_campo(t, y, p);
K2 = h * rlc_campo(t + h/2, y + K1/2, p);
K3 = h * rlc_campo(t + h/2, y + K2/2, p);
K4 = h * rlc_campo(t + h, y + K3, p);
y1 = y + (K1 + 2*K2 + 2*K3 + K4) / 6;
\end{codigo}

\paragraph{O laço de integração.} A função \texttt{integrar} recebe um dos três métodos como argumento (\texttt{passo}) e repete o passo até \(t_f\). Primeiro, ela calcula o número de passos \texttt{n} e ajusta o passo para \texttt{he}, de modo que a malha de tempo termine exatamente em \(t_f\). Dentro do laço, chama o método, soma as avaliações em \texttt{nf} e testa a divergência. Se o estado estourar, marca \texttt{estavel = \%f} e para.
\begin{codigo}
n = max(round((tf - t0) / h), 1);
he = (tf - t0) / n;
t = t0 + he * (0:n);
for k = 1:n
    [yn, m] = passo(t(k), y(:, k), he, p);
    nf = nf + m;
    if or(isnan(yn)) | or(isinf(yn)) | max(abs(yn)) > limite then
        estavel = %f;
        break
    end
    y(:, k + 1) = yn;
end
\end{codigo}

\paragraph{Soluções exatas e erro.} O arquivo \texttt{rlc\_exatas.sci} escreve as Equações \eqref{eq:i-sub} a \eqref{eq:i-res} em forma vetorial, com os operadores elemento a elemento do Scilab (\texttt{.*}). As funções \texttt{residuo\_degrau} e \texttt{residuo\_seno} conferem essas fórmulas antes do uso. A função \texttt{erro\_max} compara a solução numérica com a exata e devolve o maior erro absoluto da corrente e da tensão. Ela usa só os instantes em que as duas soluções são finitas:
\begin{codigo}
ok = and(~isnan(y), "r") & and(~isnan(ref), "r");
e = [max(abs(y(1, ok) - ref(1, ok))), max(abs(y(2, ok) - ref(2, ok)))];
\end{codigo}

\paragraph{Ordem de convergência.} A função \texttt{ordem\_observada} calcula a ordem entre dois passos vizinhos por \(\log(E_1/E_2)/\log(h_1/h_2)\). A função \texttt{ajuste\_ordem} ajusta a reta \(\log E=\log C+p\log h\) pelo método dos mínimos quadrados, resolvido com o operador \texttt{\textbackslash} do Scilab. A inclinação \(p\) é a ordem observada. A constante \(C\) alimenta a estimativa de custo da Seção~\ref{sec:resultados}.

\paragraph{Limite de estabilidade.} A função \texttt{limite\_estab} varre 40\,000 valores de \(h\), de \(0{,}125~\mu\)s a 5~ms. Para cada um, calcula \(|R(h\lambda)|\) nos dois autovalores, com a função \texttt{funcao\_estab}, que implementa a Equação~\eqref{eq:Rz}. O laço para no primeiro \(h\) instável e devolve o último \(h\) estável:
\begin{codigo}
for h = hs
    if or(abs(funcao_estab(metodo, h * p.lambda)) > 1 + 1e-9) then
        break
    end
    hmax = h;
end
\end{codigo}

\paragraph{Comparação com o comando \texttt{ode}.} A função \texttt{exp\_ode} resolve o mesmo degrau com o comando \texttt{ode} do Scilab, nas opções \texttt{"rk"} e padrão (\texttt{lsoda}), e mede o erro e o número de avaliações de \(f\) \cite[seção~10.7]{aquino2020}. Essa é a comparação com outra técnica numérica pedida no roteiro. Os métodos do trabalho não usam o \texttt{ode}.

\subsection{Como executar}

No console do Scilab, na pasta do trabalho, o comando abaixo roda todos os ensaios, grava os resultados e gera as figuras:
\begin{codigo}
exec("codigo/main.sce", -1)
\end{codigo}
Todos os números das tabelas e das figuras da Seção~\ref{sec:resultados} saem dessa execução. O programa confere as soluções exatas antes dos ensaios e para com uma mensagem de erro se algum resíduo passar de \(5\times10^{-3}\). Os resíduos medidos, da ordem de \(10^{-15}\) no degrau e de \(10^{-10}\) na senoide, mostram que o erro medido nos ensaios é do método numérico, e não da fórmula de referência.
